% Candidate method section only. Not included by the current paper.
% The original AgentDojo/ASB results have NOT been rerun with this method.
% Requires the existing paper's amsmath/amssymb packages.

\subsection{Dependency-Conditioned Joint Risk}
\label{sec:candidate-dependency-risk}

We retain DAFR's execution-boundary placement, source-aware evidence
projection, action--evidence representation, linear execution requirements,
and decision--feedback loop. The candidate replacement concerns the joint-risk
component. Instead of combining marginal risk coordinates through a norm, it
represents the conditions under which untrusted influence can produce a
specific external effect.

For the proposed action $a_t$, let $G_t$ be a source-tagged effect-dependency
graph. A path connects source influence to an effect-bearing argument and the
resulting external effect. The graph preserves the original argument field,
value, source identity, and state revision. Trusted authorization covers a
path only when its issuer, effect, field, value, and revision match the proposed
action. Let $\mathcal A_t$ contain the paths without matching authorization.
Untrusted text cannot establish authorization merely by claiming approval.

Each primitive event $X_e\in\{0,1\}$ denotes a specified condition necessary
for an unauthorized path to activate. Its meaning must be defined by the
environment model, rather than inferred from the numerical value of an
unrelated grounding or similarity score. For a path $\pi$, all its required
events must occur. The action-level unauthorized-influence event is
\begin{equation}
 U_t=\bigvee_{\pi\in\mathcal A_t}\bigwedge_{e\in\pi}(X_e=1).
 \label{eq:candidate-unauthorized-event}
\end{equation}
An empty active-path set gives $U_t=0$. The same primitive event may occur in
several paths; these paths must not be treated as independent failures.

For a declared primitive-event model with probabilities $p_{t,e}$ and
independent primitive events, inclusion--exclusion gives
\begin{equation}
 \Pr(U_t)=
 \sum_{\emptyset\ne S\subseteq\mathcal A_t}
 (-1)^{|S|+1}\prod_{e\in\bigcup_{\pi\in S}\pi}p_{t,e}.
 \label{eq:candidate-joint-probability}
\end{equation}
Taking the union of event identifiers handles dependence induced by shared
events. This assumption does not handle arbitrary dependence between distinct
primitive events. Such dependence requires a separately specified joint-event
model; the independent-event formula is not applicable by default.

Let $s_t\geq0$ be the declared consequence severity and $b_t\geq0$ the
permitted expected-unauthorized-loss budget. The joint-risk margin is
\begin{equation}
 m_t^{\mathrm{dep}}=b_t-s_t\Pr(U_t).
 \label{eq:candidate-dependency-margin}
\end{equation}
This model assumes a common action-level severity whenever any unauthorized
path activates. It does not sum separate path losses. Where consequences
differ by outcome, the loss model must instead specify those outcomes and
compute their expected loss explicitly.

The feasible set retains the individual execution requirements:
\begin{equation}
 \mathcal F_t^{\mathrm{dep}}=
 \left(\bigcap_{j\in\mathcal H_t}\mathcal P_j\right)
 \cap\{z\mid m_t^{\mathrm{dep}}(z)\geq0\}.
 \label{eq:candidate-feasible-set}
\end{equation}
Schema checks, grounding, task support, prerequisites, and authorization
requirements remain active. The joint margin does not replace these checks.
New verified evidence may change the action--evidence coordinates, graph,
active path set, or budget. Every revised proposal is checked again before
execution.

\paragraph{Uncertainty and interpretation.}
If independently justified intervals $[\ell_{t,e},u_{t,e}]$ cover the true
primitive probabilities, the monotonicity of the event in
Eq.~\eqref{eq:candidate-unauthorized-event} permits evaluation at the upper
endpoints. The resulting expected-loss bound is conditional on interval
coverage and the declared event model. LLM confidence, provenance scores,
and grounding scores are not calibrated event probabilities by themselves.
The current controlled implementation receives declared measurements and
explicit trace records; real-world probability estimation and automatic
dependency extraction remain to be established.

\paragraph{Geometry and computational cost.}
For a fixed graph and authorization context, the joint probability is a
multilinear polynomial in the primitive probabilities. Its budget sublevel
set can be nonconvex. Consequently, the original norm-region convexity
statement and general convex-region reasoning do not apply to this candidate.
The decision is evaluated directly, and feedback requests verified evidence
or scoped authorization rather than projecting to a convex feasible set.
For $K$ active paths and $n$ distinct primitive events, direct
inclusion--exclusion costs $O(2^K n)$, in addition to the linear checks and
graph construction. The current implementation evaluates small graphs; a
general linear-time or large-graph scalability claim is unsupported.

\paragraph{Scope of the candidate replacement.}
The existing AgentDojo and ASB main-table results belong to the original
halfspace-plus-risk-cone implementation. They are not evidence for this
replacement. Controlled risk-form comparisons fix evidence, scoped
authorization, uncertainty handling, verification order, and query budgets,
and replace only the joint-risk evaluator. Both a union-of-path causal world
and a quadratic-budget control world are retained, because the appropriate
risk law depends on the environment and policy semantics. Changing the main
method requires corresponding end-to-end results and component ablations.
